\section{BFS, DFS y bipartici\'on}
\whenuse{Explorar componentes, construir padres/profundidades, detectar ciclos
o colorear un grafo; BFS da distancias m\'inimas cuando todas las aristas pesan uno.}

\begin{lstlisting}[style=cpp]
vector<int> bfs_distances(const vector<vector<int>> &graph, int source) {
    int n = graph.size();
    vector<int> distance(n, -1);
    queue<int> pending;
    distance[source] = 0;
    pending.push(source);
    while (!pending.empty()) {
        int v = pending.front(); pending.pop();
        for (int to : graph[v]) if (distance[to] == -1) {
            distance[to] = distance[v] + 1;
            pending.push(to);
        }
    }
    return distance;
}

vector<int> undirected_components(const vector<vector<int>> &graph) {
    int n = graph.size(), component_count = 0;
    vector<int> component(n, -1), stack;
    for (int start = 0; start < n; ++start) if (component[start] == -1) {
        component[start] = component_count++;
        stack.push_back(start);
        while (!stack.empty()) {
            int v = stack.back(); stack.pop_back();
            for (int to : graph[v]) if (component[to] == -1) {
                component[to] = component[v];
                stack.push_back(to);
            }
        }
    }
    return component;
}

bool is_bipartite(const vector<vector<int>> &graph) {
    int n = graph.size();
    vector<int> color(n, -1);
    for (int start = 0; start < n; ++start) if (color[start] == -1) {
        queue<int> q;
        color[start] = 0; q.push(start);
        while (!q.empty()) {
            int v = q.front(); q.pop();
            for (int to : graph[v]) {
                if (color[to] == -1) color[to] = color[v] ^ 1, q.push(to);
                else if (color[to] == color[v]) return false;
            }
        }
    }
    return true;
}
\end{lstlisting}

\complexity{$O(V+E)$ tiempo y $O(V)$ memoria.}
